\documentclass[12pt]{article}
\usepackage{amsmath, amssymb, amsfonts}
\usepackage{geometry}
\usepackage{setspace}
\usepackage{enumitem}
\usepackage{booktabs}
\usepackage{graphicx}
\usepackage{hyperref}
\usepackage{siunitx}
\sisetup{
  table-number-alignment = center,
  table-figures-integer = 4,
  table-figures-decimal = 4
}

\geometry{margin=1in}
\setstretch{1.2}

\title{ECON 695 -- Problem Set 5}
\date{Due: November 3, 2025}

\begin{document}

\maketitle

\section*{AI Disclaimer}
I used AI tools (ChatGPT) to help troubleshoot minor LaTeX and Python issues, to reorganize sections for clarity (namely indentation conventions and organization in the \verb+PS5_twins.ipynb+ file), and to draft descriptive tables (because they are a pain in LaTeX). All analysis, model specification, interpretation, and final write-up are my own.

\section*{Twins Study (Princeton Twins Survey)}
In this problem set I work with the Princeton Twins Survey (\texttt{twins.csv}). Each observation is a twin from a pair of identical twins, indexed by a family id \(i\) and a twin indicator \(t\in\{1,2\}\). Below I restate the variable dictionary from the assignment for ease of reference:
\begin{itemize}[leftmargin=1.25em]
  \item \texttt{famid} = family id variable.
  \item \texttt{t} \(=1,2\) for twin \#1 or twin \#2.
  \item \texttt{age} = age (some observations include fractional years).
  \item \texttt{educ} = education (years).
  \item \texttt{oeduc} = education of the other twin.
  \item \texttt{lw} = log wage.
  \item \texttt{married} = \(1\) if married, \(0\) otherwise.
  \item \texttt{omarried} = marriage status of the other twin.
  \item \texttt{female} = \(1\) if female.
  \item \texttt{ofemale} = \(1\) if the other twin is female.
  \item \texttt{exp} = ``labor market experience'' \(=\texttt{age}-\texttt{educ}-6\).
  \item \texttt{oexp} = experience of the other twin.
\end{itemize}
\noindent\textbf{Note.} All pairs in the file are identical twins, so \(\texttt{female}=\texttt{ofemale}\) for every observation.

\subsection*{Data setup \& quick checks}
Before starting Question 1, I performed basic checks in the notebook:
\begin{enumerate}[label=\arabic*)]
  \item Uniqueness: \((\texttt{famid},\texttt{t})\) combinations are unique (one row per twin).
  \item Sex consistency: \(\texttt{female}=\texttt{ofemale}\) holds for all pairs (same sex by design).
  \item Feature engineering: I created \(\texttt{exp2}=\texttt{exp}^2\) for use in quadratic experience profiles.
\end{enumerate}

\subsection*{Descriptive statistics (key variables)}
I summarize several variables that are used repeatedly below (sample size \(n=1{,}074\)). The values match the notebook output produced via \texttt{pandas.describe()}.
\begin{table}[h!]
  \centering
  \caption{Summary statistics}
  \label{tab:desc}
  \begin{tabular}{l
                  S[table-format=4.0]
                  S[table-format=1.4]
                  S[table-format=1.4]
                  S[table-format=1.4]
                  S[table-format=1.4]
                  S[table-format=1.4]
                  S[table-format=1.4]}
    \toprule
    Variable & {count} & {mean} & {std} & {min} & {25\%} & {50\%} & {75\%} & {max} \\
    \midrule
    \texttt{lw}     & 1074 & 2.4710 & 0.6098 & 0.7226 & 2.0441 & 2.4571 & 2.8212 & 4.8929 \\
    \texttt{educ}   & 1074 & 14.0805 & 2.0298 & 8.0000 & 12.0000 & 14.0000 & 16.0000 & 18.0000 \\
    \texttt{exp}    & 1074 & 17.8336 & 11.8061 & -2.0000 & 9.0000 & 16.7361 & 25.0000 & 53.0000 \\
    \texttt{exp2}   & 1074 & 457.2928 & 518.2904 & 0.0000 & 81.0000 & 280.1114 & 625.0000 & 2809.0000 \\
    \texttt{married}& 1074 & 0.4795 & 0.4998 & 0.0000 & 0.0000 & 0.0000 & 1.0000 & 1.0000 \\
    \texttt{female} & 1074 & 0.5345 & 0.4990 & 0.0000 & 0.0000 & 1.0000 & 1.0000 & 1.0000 \\
    \bottomrule
  \end{tabular}
\end{table}

\bigskip

\section*{Q1. Simple Model: Log Wages on Schooling, Experience, Marital Status, and Sex}

\begin{quote}
Estimate a simple model relating log wages to: education, experience, experience-squared, married status, and female.
\end{quote}

\subsection*{Model}
I estimate the OLS model
\[
  \underbrace{\texttt{lw}_{it}}_{\log(\text{wage})}
  \;=\;
  \beta_0
  + \beta_1\,\texttt{educ}_{it}
  + \beta_2\,\texttt{exp}_{it}
  + \beta_3\,\texttt{exp}_{it}^{2}
  + \beta_4\,\texttt{married}_{it}
  + \beta_5\,\texttt{female}_{it}
  + \varepsilon_{it},
\]
using heteroskedasticity–robust (HC1) standard errors. The marginal effect of experience is
\( \frac{\partial\,\texttt{lw}_{it}}{\partial\,\texttt{exp}_{it}} = \beta_2 + 2\,\beta_3\,\texttt{exp}_{it} \).

\subsection*{Results}
\begin{table}[h!]
  \centering
  \caption{Baseline OLS with HC1 standard errors}
  \label{tab:q1_ols}
  \begin{tabular}{l
                  S[table-format=1.4]
                  S[table-format=1.4]
                  S[table-format=2.2]
                  S[table-format=1.3]}
    \toprule
    Variable & {Coefficient} & {Robust SE} & {$t$} & {$p$} \\
    \midrule
    Intercept       & 0.3455 & 0.1386 &  2.49 & 0.013 \\
    Education       & 0.1238 & 0.0086 & 14.36 & 0.000 \\
    Experience      & 0.0523 & 0.0048 & 10.88 & 0.000 \\
    Experience$^{2}$& -0.0009& 0.0001 & -8.28 & 0.000 \\
    Married         & 0.0693 & 0.0321 &  2.16 & 0.031 \\
    Female          & -0.3341& 0.0309 & -10.83& 0.000 \\
    \midrule
    \multicolumn{5}{l}{\footnotesize $N=1{,}074$;\quad $R^2=0.3146$;\quad $\bar R^{2}=0.3113$.} \\
    \bottomrule
  \end{tabular}
\end{table}

\subsection*{Interpretation}
\begin{itemize}[leftmargin=1.25em]
  \item Schooling: an additional year of education raises expected log wages by about \(0.124\) log points; using the exact transformation, this is \(100\times(\exp(0.1238)-1)\approx 13.2\%\).
  \item Experience: wages rise with experience but at a diminishing rate because \(\beta_3<0\). The implied turning point is
  \[
    \texttt{exp}^{*} \;=\; \frac{\beta_2}{-2\beta_3}
    \;=\; \frac{0.0523}{2(0.0009)} \;\approx\; 29.1 \text{ years},
  \]
  beyond which additional experience adds little to wages on average.
  \item Marital status: being married is associated with roughly \(100\times(\exp(0.0693)-1)\approx 7.2\%\) higher wages, significant at the \(5\%\) level.
  \item Sex: the coefficient on \(\texttt{female}\) is \(-0.334\); in levels this corresponds to \(100\times(\exp(-0.3341)-1)\approx -28.4\%\), i.e., women earn about \(28\%\) less than men on average, conditional on covariates.
  \item Model fit: the model explains about \(31\%\) of the variation in \(\log(\text{wage})\) (\(R^2=0.3146\)), which is reasonable for a cross-sectional wage regression.
\end{itemize}


\section*{Q2. Separate Models for Men and Women}

\begin{quote}
Fit the same model separately for men and women. Does the model look different?
\end{quote}

\subsection*{Specification}
For each gender \(g\in\{\text{men},\text{women}\}\) I estimate
\[
  \texttt{lw}_{i} \;=\; \beta_{0g}
  + \beta_{1g}\,\texttt{educ}_{i}
  + \beta_{2g}\,\texttt{exp}_{i}
  + \beta_{3g}\,\texttt{exp}_{i}^{2}
  + \beta_{4g}\,\texttt{married}_{i}
  + u_{i},
\]
using HC1 robust standard errors.

\subsection*{Results}

% --- Q2 results: separate, non-overlapping tables ---
\begin{table}[htbp]
  \centering
  \caption{Men: OLS with HC1 robust standard errors}
  \label{tab:q2_men}
  \small
  \begin{tabular}{l
                  S[table-format=1.4]
                  S[table-format=1.4]
                  S[table-format=2.4]
                  S[table-format=1.4]}
    \toprule
    Variable & {Coef.} & {SE(HC1)} & {$t$} & {$p$} \\
    \midrule
    Intercept        & 0.5197 & 0.2194 &  2.3685 & 0.0179 \\
    Education        & 0.1027 & 0.0135 &  7.5896 & 0.0000 \\
    Experience       & 0.0588 & 0.0080 &  7.3951 & 0.0000 \\
    Experience$^{2}$ & -0.0010& 0.0002 & -5.8950 & 0.0000 \\
    Married          & 0.1993 & 0.0506 &  3.9419 & 0.0001 \\
    \midrule
    \multicolumn{5}{l}{\footnotesize $N=500$;\; $R^2=0.2553$;\; $\bar R^{2}=0.2493$.} \\
    \bottomrule
  \end{tabular}
\end{table}

\begin{table}[htbp]
  \centering
  \caption{Women: OLS with HC1 robust standard errors}
  \label{tab:q2_women}
  \small
  \begin{tabular}{l
                  S[table-format=1.4]
                  S[table-format=1.4]
                  S[table-format=2.4]
                  S[table-format=1.4]}
    \toprule
    Variable & {Coef.} & {SE(HC1)} & {$t$} & {$p$} \\
    \midrule
    Intercept        & -0.0854 & 0.1595 & -0.5353 & 0.5924 \\
    Education        & 0.1369  & 0.0101 & 13.5113 & 0.0000 \\
    Experience       & 0.0529  & 0.0059 &  8.9736 & 0.0000 \\
    Experience$^{2}$ & -0.0010 & 0.0001 & -6.8237 & 0.0000 \\
    Married          & -0.0399 & 0.0407 & -0.9813 & 0.3264 \\
    \midrule
    \multicolumn{5}{l}{\footnotesize $N=574$;\; $R^2=0.3033$;\; $\bar R^{2}=0.2984$.} \\
    \bottomrule
  \end{tabular}
\end{table}

\subsection*{Turning points}
With the quadratic experience term, the implied turning point is
\[
  \texttt{exp}^{*}_g \;=\; \frac{\beta_{2g}}{-2\,\beta_{3g}}.
\]
Plugging in the estimates gives
\[
  \text{Men: } \texttt{exp}^{*}\approx 29.8 \text{ years}
  \qquad\text{and}\qquad
  \text{Women: } \texttt{exp}^{*}\approx 26.9 \text{ years}.
\]

\subsection*{Does the model look different?}
\begin{itemize}[leftmargin=1.25em]
  \item \textbf{Education:} Positive and highly significant for both; the effect is \emph{larger for women} (\(0.137\) vs. \(0.103\)).
  \item \textbf{Experience profile:} Positive \(\beta_2\) and negative \(\beta_3\) in both groups (concave), with a slightly earlier peak for women (26.9 vs.\ 29.8 years).
  \item \textbf{Married:} \emph{Men:} positive and significant (\(\approx 0.20\)). \emph{Women:} small, negative, and insignificant (\(-0.040\)).
  \item \textbf{Fit:} \(R^2\) is a bit higher for women (0.303) than men (0.255).
\end{itemize}

\section*{Q3. Mean Family Marriage Rate}

\begin{quote}
Construct the mean family marriage rate for each person in the data set (i.e., the fraction of the twins that is married, which can be 0, 1/2, or 1).
\end{quote}

\subsection*{Part (a): Verifying the Auxiliary Regression}

We define the mean family marriage rate for each family \(f\) as the average of the twins’ marriage indicators:
\[
  \texttt{mean\_married}_f \;=\; \frac{1}{2}\sum_{t=1}^{2}\texttt{married}_{ft},
\]
which can take values \(0\), \(0.5\), or \(1\), depending on whether neither, one, or both twins are married.

We then estimate the auxiliary regression
\[
  \texttt{married}_{ft} = \alpha + \rho\,\texttt{mean\_married}_{f} + \varepsilon_{ft},
\]
using HC1 robust standard errors.

\begin{table}[h!]
  \centering
  \caption{Auxiliary regression: verifying the mean–marriage relationship}
  \label{tab:q3a_aux}
  \small
  \begin{tabular}{l
                  S[table-format=1.4]
                  S[table-format=1.4]
                  S[table-format=2.2]
                  S[table-format=1.3]}
    \toprule
    Variable & {Coefficient} & {Robust SE (HC1)} & {$t$} & {$p$} \\
    \midrule
    Intercept       & -0.0000 & 0.0074 & -0.0000 & 1.0000 \\
    Mean Married    &  1.0000 & 0.0009 & 1147.51 & 0.0000 \\
    \midrule
    \multicolumn{5}{l}{\footnotesize $N=1{,}074$;\quad $R^2=0.7370$;\quad $\bar R^{2}=0.7367$.} \\
    \bottomrule
  \end{tabular}
\end{table}

\subsubsection*{Interpretation}
The estimated coefficient on \(\texttt{mean\_married}\) is essentially 1, confirming the theoretical result discussed in Lecture 8.  
Intuitively, each twin’s marital status is mechanically tied to the family’s mean marriage rate: if one twin’s status changes (from 0 to 1, or vice versa), the family average changes in the same direction by exactly half for both siblings, implying a perfect one-to-one mapping in the regression. 

Mathematically, for any twin \(t\) in family \(f\),
\[
  \texttt{mean\_married}_{f}
  = \frac{1}{2}\big(\texttt{married}_{f1} + \texttt{married}_{f2}\big)
  = \frac{1}{2}\texttt{married}_{ft} + \frac{1}{2}\texttt{married}_{f(-t)}.
\]
Because \(\texttt{mean\_married}_f\) contains \(\texttt{married}_{ft}\) as part of its own definition, regressing \(\texttt{married}_{ft}\) on \(\texttt{mean\_married}_f\) yields an estimated slope of 1 by construction: any within-family change in \(\texttt{married}_{ft}\) moves both sides of the equation equally.

Empirically, this is reflected in the near-perfect fit (\(R^2\approx0.74\)) and a coefficient statistically indistinguishable from 1.

\subsection*{Part (b): Adding the family mean marriage rate to gender-specific wage models}

We augment the Q2 models to include the family mean marriage rate \(\texttt{mean\_married}_f\) and estimate, separately by gender,
\[
  \texttt{lw}_{i} \;=\; \beta_0 + \beta_1\,\texttt{educ}_{i}
  + \beta_2\,\texttt{exp}_{i} + \beta_3\,\texttt{exp}_{i}^{2}
  + \beta_4\,\texttt{married}_{i}
  + \beta_5\,\texttt{mean\_married}_{f} + u_{i},
\]
using HC1 robust standard errors.

% --- Q3(b) results: Men ---
\begin{table}[htbp]
  \centering
  \caption{Men: OLS with \texttt{mean\_married} (HC1 robust)}
  \label{tab:q3b_men}
  \small
  \begin{tabular}{l
                  S[table-format=1.4]
                  S[table-format=1.4]
                  S[table-format=2.4]
                  S[table-format=1.4]}
    \toprule
    Variable & {Coef.} & {SE(HC1)} & {$t$} & {$p$} \\
    \midrule
    Intercept         & 0.5102 & 0.2200 &  2.3187 & 0.0204 \\
    Education         & 0.1029 & 0.0135 &  7.6062 & 0.0000 \\
    Experience        & 0.0572 & 0.0079 &  7.1902 & 0.0000 \\
    Experience$^{2}$  & -0.0010& 0.0002 & -5.8788 & 0.0000 \\
    Married           & 0.0300 & 0.0897 &  0.3339 & 0.7385 \\
    Mean married      & 0.2500 & 0.1108 &  2.2571 & 0.0240 \\
    \midrule
    \multicolumn{5}{l}{\footnotesize $N=500$;\; $R^2=0.2619$;\; $\bar R^{2}=0.2544$.} \\
    \bottomrule
  \end{tabular}
\end{table}

% --- Q3(b) results: Women ---
\begin{table}[htbp]
  \centering
  \caption{Women: OLS with \texttt{mean\_married} (HC1 robust)}
  \label{tab:q3b_women}
  \small
  \begin{tabular}{l
                  S[table-format=1.4]
                  S[table-format=1.4]
                  S[table-format=2.4]
                  S[table-format=1.4]}
    \toprule
    Variable & {Coef.} & {SE(HC1)} & {$t$} & {$p$} \\
    \midrule
    Intercept         & -0.0826 & 0.1597 & -0.5173 & 0.6050 \\
    Education         & 0.1368  & 0.0102 & 13.4580 & 0.0000 \\
    Experience        & 0.0534  & 0.0062 &  8.6071 & 0.0000 \\
    Experience$^{2}$  & -0.0010 & 0.0001 & -6.6896 & 0.0000 \\
    Married           & -0.0221 & 0.0728 & -0.3031 & 0.7618 \\
    Mean married      & -0.0268 & 0.0915 & -0.2930 & 0.7695 \\
    \midrule
    \multicolumn{5}{l}{\footnotesize $N=574$;\; $R^2=0.3034$;\; $\bar R^{2}=0.2972$.} \\
    \bottomrule
  \end{tabular}
\end{table}

% --- compact comparison of marriage coefficients ---
\begin{table}[htbp]
  \centering
  \caption{Marriage coefficients before (Q2) vs after adding \texttt{mean\_married} (Q3b)}
  \label{tab:q3b_compare}
  \small
  \begin{tabular}{l
                  S[table-format=1.4]
                  S[table-format=1.4]
                  S[table-format=1.4]
                  S[table-format=+1.4]}
    \toprule
    Group & {Married (Q2)} & {Married (Q3b)} & {Mean married} & {$\Delta$ Married (Q3b$-$Q2)} \\
    \midrule
    Men    & 0.1993 & 0.0300 & 0.2500 & -0.1693 \\
    Women  & -0.0399 & -0.0221 & -0.0268 & 0.0179 \\
    \bottomrule
  \end{tabular}
\end{table}

\subsubsection*{What changed?}
\begin{itemize}[leftmargin=1.25em]
  \item \textbf{Men.} The individual marriage coefficient shrinks from \(\approx 0.199\) (Q2) to \(\approx 0.030\) (Q3b) and becomes insignificant, while the family mean term is positive and significant (\(\widehat\beta_5\approx 0.25\)).
  \item \textbf{Women.} The marriage coefficient was already small and insignificant in Q2 (\(\approx -0.040\)) and remains small/insignificant in Q3b; the family mean term is also small and insignificant.
\end{itemize}

\subsubsection*{Interpretation (men vs. women)}
With two siblings per family, \(\texttt{mean\_married}_f=(\texttt{married}_{f1}+\texttt{married}_{f2})/2\) captures a \emph{family-level (shared) component} of marriage—selection or common factors that influence both twins (values, timing, labor-market conditions).  
For men, most of the apparent “marriage premium” in Q2 is absorbed by this shared component (strong \(\widehat\beta_5>0\)), leaving little individual-specific premium after controlling for the family mean.  
For women, there is no evident premium to begin with, and controlling for the family mean changes little—both the individual and family components are economically small and statistically insignificant.

\subsubsection*{Model fit}
Fit is essentially unchanged for women (\(R^2 \approx 0.303 \to 0.303\)) and rises only slightly for men (\(0.255 \to 0.262\)), consistent with explanatory power from the family mean operating mainly in the male sample.

\subsection*{Part (c): Using the co-twin's marriage status instead of the family mean}

We re-estimate the gender-specific models from Q2 replacing the family mean term with the
other twin’s marriage status \(\texttt{omarried}_{i}\):
\[
  \texttt{lw}_{i}
  \;=\; \beta_0 + \beta_1\,\texttt{educ}_{i}
  + \beta_2\,\texttt{exp}_{i} + \beta_3\,\texttt{exp}_{i}^{2}
  + \beta_4\,\texttt{married}_{i} + \beta_5\,\texttt{omarried}_{i} + u_{i},
\]
using HC1 robust SEs.

% --- Q3(c) results: Men ---
\begin{table}[htbp]
  \centering
  \caption{Men: OLS with \texttt{omarried} (HC1 robust)}
  \label{tab:q3c_men}
  \small
  \begin{tabular}{l
                  S[table-format=1.4]
                  S[table-format=1.4]
                  S[table-format=2.4]
                  S[table-format=1.4]}
    \toprule
    Variable & {Coef.} & {SE(HC1)} & {$t$} & {$p$} \\
    \midrule
    Intercept        & 0.5102 & 0.2200 &  2.3187 & 0.0204 \\
    Education        & 0.1029 & 0.0135 &  7.6062 & 0.0000 \\
    Experience       & 0.0572 & 0.0079 &  7.1902 & 0.0000 \\
    Experience$^{2}$ & -0.0010& 0.0002 & -5.8788 & 0.0000 \\
    Married          & 0.1550 & 0.0539 &  2.8771 & 0.0040 \\
    Other twin married & 0.1250 & 0.0554 &  2.2571 & 0.0240 \\
    \midrule
    \multicolumn{5}{l}{\footnotesize $N=500$;\; $R^2=0.2619$;\; $\bar R^{2}=0.2544$.} \\
    \bottomrule
  \end{tabular}
\end{table}

% --- Q3(c) results: Women ---
\begin{table}[htbp]
  \centering
  \caption{Women: OLS with \texttt{omarried} (HC1 robust)}
  \label{tab:q3c_women}
  \small
  \begin{tabular}{l
                  S[table-format=1.4]
                  S[table-format=1.4]
                  S[table-format=2.4]
                  S[table-format=1.4]}
    \toprule
    Variable & {Coef.} & {SE(HC1)} & {$t$} & {$p$} \\
    \midrule
    Intercept        & -0.0826 & 0.1597 & -0.5173 & 0.6050 \\
    Education        & 0.1368  & 0.0102 & 13.4580 & 0.0000 \\
    Experience       & 0.0534  & 0.0062 &  8.6071 & 0.0000 \\
    Experience$^{2}$ & -0.0010 & 0.0001 & -6.6896 & 0.0000 \\
    Married          & -0.0355 & 0.0432 & -0.8202 & 0.4121 \\
    Other twin married & -0.0134 & 0.0457 & -0.2930 & 0.7695 \\
    \midrule
    \multicolumn{5}{l}{\footnotesize $N=574$;\; $R^2=0.3034$;\; $\bar R^{2}=0.2972$.} \\
    \bottomrule
  \end{tabular}
\end{table}

% --- concise comparison with Q2 and Q3(b) ---
\begin{table}[htbp]
  \centering
  \caption{Marriage coefficients across specifications}
  \label{tab:q3c_compare}
  \small
  \begin{tabular}{l
                  S[table-format=+1.4]
                  S[table-format=+1.4]
                  S[table-format=+1.4]
                  S[table-format=+1.4]
                  S[table-format=+1.4]
                  S[table-format=+1.4]
                 }
    \toprule
    Group & {Married (Q2)} & {Married (Q3b)} & {Mean married (Q3b)} & {Married (Q3c)} & {Other twin (Q3c)} \\
    \midrule
    Men   &  0.1993 &  0.0300 &  0.2500 &  0.1550 &  0.1250 \\
    Women & -0.0399 & -0.0221 & -0.0268 & -0.0355 & -0.0134 \\
    \bottomrule
  \end{tabular}
\end{table}

\subsubsection*{Comparison and interpretation}
\paragraph{Link between (b) and (c).}
With two siblings per family,
\[
  \texttt{mean\_married}_f
  = \tfrac12\big(\texttt{married}_{f1}+\texttt{married}_{f2}\big)
  = \tfrac12\,\texttt{married}_{i} + \tfrac12\,\texttt{omarried}_{i}.
\]
In a linear model, this implies the \emph{mean-marriage} slope in (b) should be close to the average of the two slopes in (c):
\[
  \widehat\beta^{(b)}_{\text{mean}} \;\approx\; \tfrac12
  \big(\widehat\beta^{(c)}_{\text{married}}+\widehat\beta^{(c)}_{\text{omarried}}\big).
\]
Our estimates line up: for men,
\(\widehat\beta^{(c)}_{\text{married}}+\widehat\beta^{(c)}_{\text{omarried}}
=0.155+0.125=0.280\), so \(\tfrac12(\cdot)\approx 0.140\), which is close in magnitude to the strong family-mean effect in (b) (\(\approx0.25\)) and explains why adding \(\texttt{mean\_married}\) in (b) pulled the own-marriage coefficient toward zero.

\paragraph{Men.}
Relative to Q2, the individual marriage premium falls when adding \(\texttt{mean\_married}\) (Q3b), but when we instead add \(\texttt{omarried}\) (Q3c), both own marriage (\(0.155\), \(p=0.004\)) and co-twin marriage (\(0.125\), \(p=0.024\)) are positive and significant. Taken together, they recover a combined effect near the family-mean channel from (b). This pattern is consistent with a shared (family/selection) component to the “marriage premium” among men rather than a large individual-specific causal effect.

\paragraph{Women.}
Coefficients on both \(\texttt{married}\) and \(\texttt{omarried}\) remain small and statistically insignificant, echoing Q2–Q3(b). Controlling for the family mean or the co-twin’s status does not reveal a meaningful association between marriage and wages for women in this sample.

\paragraph{Bottom line.}
Replacing \(\texttt{mean\_married}\) with \(\texttt{omarried}\) leaves the substantive story intact: the apparent premium in men is largely accounted for by shared (family-level) factors that influence both twins; for women, marriage has no robust association with wages across specifications.

\section*{Q4. Three Ways to Get the Within Estimator (men only)}

\subsection*{(a) Within estimator via family demeaning}

\paragraph{Setup.}
Restrict the sample to \emph{male} twins ($N=500$). For each family $f$ with twins $t\in\{1,2\}$,
apply the within (demeaning) transformation
\[
  \tilde y_{ft} \equiv y_{ft}-\bar y_f, \qquad
  \tilde x_{ft} \equiv x_{ft}-\bar x_f ,
\]
where $x_{ft}\in\{\texttt{educ},\,\texttt{exp}^2,\,\texttt{married}\}$ and $y_{ft}=\texttt{lw}_{ft}$.
Because experience is defined as $\texttt{exp}=\texttt{age}-\texttt{educ}-6$ and identical twins are the
same age, $\widetilde{\texttt{exp}}_{ft} = -\,\widetilde{\texttt{educ}}_{ft}$ within family. Hence
$\tilde{\texttt{exp}}$ is perfectly collinear with $\tilde{\texttt{educ}}$ (corr.\ $=-1$), so we drop
$\tilde{\texttt{exp}}$ and keep $\widetilde{\texttt{exp}^2}$.

\paragraph{Model.}
We estimate the within regression (no intercept after demeaning):
\[
  \tilde{\texttt{lw}}_{ft}
  \;=\;
  \gamma_1\,\tilde{\texttt{educ}}_{ft}
  \;+\;
  \gamma_3\,\widetilde{\texttt{exp}^2}_{ft}
  \;+\;
  \gamma_4\,\tilde{\texttt{married}}_{ft}
  \;+\; u^{W}_{ft},
\]
with HC1 robust standard errors.

\begin{table}[htbp]
  \centering
  \caption{Within-family (demeaned) regression for men (HC1 robust)}
  \label{tab:q4a_within_men}
  \small
  \begin{tabular}{l
                  S[table-format=1.4]
                  S[table-format=1.4]
                  S[table-format=2.4]
                  S[table-format=1.4]}
    \toprule
    Variable & {Coef.} & {SE(HC1)} & {$t$} & {$p$} \\
    \midrule
    $\tilde{\texttt{educ}}$         & 0.0851 & 0.0369 &  2.3087 & 0.0210 \\
    $\widetilde{\texttt{exp}^2}$    & 0.0004 & 0.0010 &  0.4402 & 0.6598 \\
    $\tilde{\texttt{married}}$      & 0.0247 & 0.0470 &  0.5261 & 0.5988 \\
    \midrule
    \multicolumn{5}{l}{\footnotesize $N=500$;\; $R^2=0.0357$.} \\
    \bottomrule
  \end{tabular}
\end{table}

\subsubsection*{Interpretation}
\begin{itemize}[leftmargin=1.25em]
  \item The within estimator compares twins \emph{within the same family}, removing all time-invariant shared family factors.  
  \item A one-year within-pair difference in schooling is associated with about \(0.085\) higher log wages for the more educated brother (significant at 5\%).  
  \item The coefficients on the demeaned quadratic experience term and marital status are small and imprecise, suggesting little additional \emph{within-pair} wage difference from these factors once shared family effects are removed.  
\end{itemize} 


\subsection*{(b) Fixed Effects}

\paragraph{Model.}
We now estimate the wage equation for male twins using a \emph{family fixed effects} model, which controls for all unobserved, time-invariant family characteristics shared by the twins (e.g., parental background, neighborhood). The regression is
\[
  \texttt{lw}_{ft} \;=\; \alpha_f \;+\; \beta_1\,\texttt{educ}_{ft}
  \;+\; \beta_3\,\texttt{exp}^2_{ft}
  \;+\; \beta_4\,\texttt{married}_{ft} \;+\; u_{ft},
\]
where $\alpha_f$ denotes the family fixed effect. As in part (a), we omit $\texttt{exp}_{ft}$ because, for identical twins of the same age, $\widetilde{\texttt{exp}}_{ft}$ is perfectly collinear with $\widetilde{\texttt{educ}}_{ft}$ when conditioning on $\alpha_f$.

\begin{table}[htbp]
  \centering
  \caption{Fixed effects (men): family FE with HC1 robust SEs}
  \label{tab:q4b_fe_men}
  \small
  \begin{tabular}{l
                  S[table-format=1.4]
                  S[table-format=1.4]
                  S[table-format=2.4]
                  S[table-format=1.4]}
    \toprule
    Variable & {Coef.} & {SE(HC1)} & {$t$} & {$p$} \\
    \midrule
    Education         & 0.0851 & 0.0523 & 1.6276 & 0.1049 \\
    Experience$^{2}$  & 0.0004 & 0.0014 & 0.3103 & 0.7566 \\
    Married           & 0.0247 & 0.0666 & 0.3709 & 0.7110 \\
    \midrule
    \multicolumn{5}{l}{\footnotesize $N=500$;\; $R^2_{\text{within}}=0.0357$;\; $R^2_{\text{overall}}=0.7490$.} \\
    \bottomrule
  \end{tabular}
\end{table}

\subsubsection*{Interpretation}
\begin{itemize}[leftmargin=1.25em]
  \item The family FE specification compares twins within the same family while allowing each family to have its own intercept $\alpha_f$. This removes all time-invariant shared background factors.
  \item The schooling coefficient (\(\widehat\beta_1\approx0.085\)) is positive but not statistically significant at conventional levels (\(p\approx0.105\)); economically, it is close to the within (demeaned) estimate from part (a), indicating both approaches recover the same within estimator.
  \item The coefficients on \(\texttt{exp}^2\) and \(\texttt{married}\) are small and imprecise, suggesting limited additional \emph{within-family} wage differences associated with these covariates once shared family effects are absorbed.
  \item The reported \(R^2_{\text{within}}\) (\(\approx0.036\)) summarizes variation explained \emph{within families}; the much larger \(R^2_{\text{overall}}\) (\(\approx0.749\)) reflects that most cross-family wage differences are captured by the family intercepts \(\alpha_f\).
\end{itemize}

\subsection*{(c) Control Function (within by adding family means)}

\paragraph{Model.}
We estimate the \emph{control function} version of the within estimator for \textit{male} twins by regressing log wages on the individual covariates and their family means:
\[
  \texttt{lw}_{ft}
  \;=\; \beta_1\,\texttt{educ}_{ft}
  \;+\; \beta_3\,\texttt{exp}^2_{ft}
  \;+\; \beta_4\,\texttt{married}_{ft}
  \;+\; \delta_1\,\bar{\texttt{educ}}_{f}
  \;+\; \delta_3\,\overline{\texttt{exp}^2}_{f}
  \;+\; \delta_4\,\overline{\texttt{married}}_{f}
  \;+\; u_{ft},
\]
using HC1 robust standard errors. Writing \(x_{ft}=(x_{ft}-\bar x_f)+\bar x_f\) shows that including \(\bar x_f\) absorbs shared family background; the coefficients on the \emph{individual} regressors
\((\beta_1,\beta_3,\beta_4)\) coincide with the slopes from the demeaned and family–fixed–effects regressions.

\begin{table}[htbp]
  \centering
  \caption{Control function (men): individual covariates and family means (HC1 robust)}
  \label{tab:q4c_ctrl_men}
  \small
  \begin{tabular}{l
                  S[table-format=1.4]
                  S[table-format=1.4]
                  S[table-format=2.4]
                  S[table-format=1.4]}
    \toprule
    Variable & {Coef.} & {SE(HC1)} & {$t$} & {$p$} \\
    \midrule
    Education                     & 0.0851 & 0.0689 & 1.2352 & 0.2168 \\
    Experience$^{2}$              & 0.0004 & 0.0018 & 0.2366 & 0.8130 \\
    Married                       & 0.0247 & 0.0987 & 0.2504 & 0.8023 \\
    Family mean education         & 0.0030 & 0.0700 & 0.0431 & 0.9656 \\
    Family mean experience$^{2}$  & -0.0003& 0.0018 & -0.1426& 0.8866 \\
    Family mean married           & 0.3590 & 0.1199 & 2.9934 & 0.0028 \\
    \midrule
    \multicolumn{5}{l}{\footnotesize $N=500$;\; $R^2=0.1631$.} \\
    \bottomrule
  \end{tabular}
\end{table}

\paragraph{Equivalence check.}
The estimated slopes on the \emph{individual} regressors match those from parts (a) and (b):

\begin{table}[h!]
  \centering
  \caption{Slopes on individual covariates are identical across (a), (b), and (c)}
  \label{tab:q4c_compare}
  \small
  \begin{tabular}{l
                  S[table-format=1.4]
                  S[table-format=1.4]
                  S[table-format=1.4]}
    \toprule
    Variable & {(a) Demeaned} & {(b) Family FE} & {(c) Control fn.} \\
    \midrule
    Education         & 0.0851 & 0.0851 & 0.0851 \\
    Experience$^{2}$  & 0.0004 & 0.0004 & 0.0004 \\
    Married           & 0.0247 & 0.0247 & 0.0247 \\
    \bottomrule
  \end{tabular}
\end{table}

\subsubsection*{Interpretation and algebraic link}
Because identical twins share age, \(\widetilde{\texttt{exp}}_{ft}=-\,\widetilde{\texttt{educ}}_{ft}\), so \(\texttt{exp}\) drops from the within/FE equations and only \(\texttt{exp}^2\) varies within family. Decompose each covariate:
\[
  x_{ft}=(x_{ft}-\bar x_f)+\bar x_f.
\]
Substituting into a linear model yields
\[
  \texttt{lw}_{ft}
  = \underbrace{\beta^{\mathrm W}}_{\text{within slope}}(x_{ft}-\bar x_f)
  + \underbrace{\delta}_{\text{between slope}}\bar x_f + \varepsilon_{ft}.
\]
Therefore, the coefficient on \(x_{ft}\) in \emph{(c)} equals the within slope \(\beta^{\mathrm W}\) recovered in \emph{(a)} (demeaning) and \emph{(b)} (family FE). Numerically, our estimates confirm this equivalence (Tables \ref{tab:q4a_within_men}, \ref{tab:q4b_fe_men}, \ref{tab:q4c_ctrl_men}, \ref{tab:q4c_compare}). Economically, the within-family schooling effect is about \(0.085\) log points per year, while within-family variation in \(\texttt{exp}^2\) and \(\texttt{married}\) explains little. The positive, significant family-mean coefficient on \(\overline{\texttt{married}}_f\) captures a \emph{shared} marriage component (selection/common factors) rather than an individual-specific premium.

\end{document}
